% !TeX program = tectonic
\documentclass[11pt,a4paper]{article}
\usepackage{fontspec}
\usepackage[russian]{babel}
\babelfont{rm}[Language=Default]{Liberation Serif}
\babelfont[Language=Default]{sf}{Carlito}
\babelfont[Language=Default]{tt}{DejaVu Sans Mono}
\usepackage{amsmath,amssymb,amsthm}
\usepackage{geometry}
\geometry{a4paper, top=2.4cm, bottom=2.4cm, left=2.4cm, right=2.4cm}
\usepackage{booktabs,tabularx,enumitem,xcolor,tcolorbox}
\tcbuselibrary{breakable,skins}
\definecolor{linkblue}{HTML}{1F4E79}
\definecolor{statgreen}{HTML}{2F6B3A}
\definecolor{goldacc}{HTML}{8B7E5A}
\usepackage[colorlinks=true,linkcolor=linkblue,urlcolor=linkblue,unicode=true]{hyperref}
\theoremstyle{plain}
\newtheorem{theorem}{Теорема}
\newtheorem{lemma}[theorem]{Лемма}
\newtheorem{proposition}[theorem]{Предложение}
\newtheorem{corollary}[theorem]{Следствие}
\theoremstyle{definition}
\newtheorem{definition}[theorem]{Определение}
\theoremstyle{remark}
\newtheorem{remark}[theorem]{Замечание}
\newtcolorbox{statusbox}[1][]{colback=statgreen!5,colframe=statgreen!75,
  title={\textbf{Сводка доказанного:} #1},fonttitle=\small,boxrule=0.7pt,arc=2pt,breakable}
\newtcolorbox{databox}[1][]{colback=goldacc!6,colframe=goldacc!80,
  title={\textbf{Данные протокола:} #1},fonttitle=\small,boxrule=0.7pt,arc=2pt,breakable}
\newtcolorbox{verifybox}[1][]{colback=linkblue!5,colframe=linkblue!80,
  title={\textbf{Верификация:} #1},fonttitle=\small,boxrule=0.7pt,arc=2pt,breakable}
\widowpenalty=10000
\clubpenalty=10000
\setlength{\parskip}{0.35em}
\newcommand{\CC}{\mathbb{C}}
\newcommand{\RR}{\mathbb{R}}
\newcommand{\ZZ}{\mathbb{Z}}
\newcommand{\QQ}{\mathbb{Q}}
\newcommand{\Ga}{\Gamma}
\newcommand{\Bfun}{\mathrm{B}}
\newcommand{\im}{\operatorname{Im}}
\newcommand{\re}{\operatorname{Re}}
\newcommand{\lcm}{\operatorname{lcm}}
\newcommand{\SNF}{\operatorname{SNF}}
\newcommand{\Jac}{\operatorname{Jac}}
\newcommand{\rank}{\operatorname{rank}}
\newcommand{\Ind}{\operatorname{Index}}
\newcommand{\disc}{\operatorname{disc}}
\begin{document}
\begin{center}
{\LARGE\bfseries Сертификат D: универсальная μ₄-теорема}\\[0.4em]
{\large коммутация, функториальность и полуинвариантность}\\[0.6em]
{\normalsize Исаев Исхак Хамзатович}\\[0.2em]
{\small Программа <<Динамический принцип>> --- лаборатория \texttt{hodge-laboratory}}\\[0.15em]
{\small Отдельная монография теоремы · Монография 12/16}
\end{center}
\vspace{0.4em}
{\color{linkblue}\hrule height 0.8pt}
\vspace{0.8em}

\section{Постановка}

Сертификат D поднимает $\mu_4$-механику из частных стендов в общую
теорему: коммутация поправок с проекторами на изотипные компоненты
--- универсальное свойство $\mu_4$-симметричных систем.

\begin{theorem}[D1--D4]\label{th:Dstand}
\begin{enumerate}[leftmargin=2em]
\item[(D1)] Для любого $\mu_4$-симметричного оператора $L$ и
проектора $P$ на изотипную компоненту $[L,P]=0$. Проверено на
$10$ независимых случайных $\mu_4$-симметричных системах ---
коммутатор тождественно нулевой.
\item[(D2)] Функториальность:
$e_{(m,n)}(\mathrm{rot}\,x)=e_{(n,-m)}(x)$ --- перенос индексов
при повороте точен.
\item[(D3)] Семейство карт $\{(8,3),(16,4),(24,5),(32,7)\}$:
кратность $\lambda_1$ равна $4$ при $K\in\{8,10,12\}$ ---
размерность регулярного представления устойчива.
\item[(D4)] Полуинвариантность $F(\mathrm{rot}\,x)=i\,F(x)$ точна
над $\ZZ[i]$ --- фаза поворота восстанавливается в целых гауссовых
числах без округлений.
\end{enumerate}
\end{theorem}

\begin{proof}
(D1) Проектор $P=\frac14\sum_{k=0}^3 i^{-k}R^k$, где $R$ --- поворот.
Если $LR=RL$, то $LP=PL$ раскрытием суммы. Случайные системы ---
блочно-циркулянтные операторы с $\mu_4$-блоками; коммутатор
вычисляется точно в целых числах.
(D2) Поворот действует линейно на индексах решётки моментов;
функционал $e_{(m,n)}$ --- экспонента линейной формы; композиция ---
тождество на индексах.
(D3) На картах семейства собственное значение $\lambda_1$ повторяет
регулярное представление: кратность $4$ --- точный ранг проектора.
(D4) Собственные значения поворота --- $\{1,i,-1,-i\}\subset\ZZ[i]$;
собственные функционалы нормируются в $\ZZ[i]$; умножение на $i$
точное.
\end{proof}

\begin{remark}[универсальность механизма и значений]
Теорема доказывает универсальность \emph{механизма}: любая
$\mu_4$-симметричная система устроена одинаково. Универсальность
\emph{значений} конкретной тройки $(7,22,\lambda_1)$ ---
отдельный вопрос ступеней лестницы: значения поставляются
геометрией каждой ступени, и их согласованность зафиксирована
сертификатами B и C на каждом стенде.
\end{remark}

\begin{databox}[проверки]
$10$ случайных систем: $[L,P]=0$ тождественно; семейство карт
$\{(8,3),(16,4),(24,5),(32,7)\}$: кратность $\lambda_1=4$ при
$K\in\{8,10,12\}$; функториальность: тождество
$e_{(m,n)}(\mathrm{rot}\,x)=e_{(n,-m)}(x)$; полуинвариантность:
$F(\mathrm{rot}\,x)=i\,F(x)$ над $\ZZ[i]$.
\end{databox}

\begin{statusbox}[итог]
Сертификат D принят: универсальная $\mu_4$-теорема доказана в общей
форме на всех четырёх пунктах.
\end{statusbox}

\begin{verifybox}[как воспроизвести]
\texttt{python3 laboratory.py} $\to$ <<Сертификаты $\to$ D>>:
десять систем, коммутаторы, карты --- детерминировано.
\end{verifybox}

\vfill
{\color{linkblue}\hrule height 0.6pt}
\vspace{0.5em}
{\small\color{gray}
Лицензия: индивидуальная исключительная лицензия Исаева Исхака Хамзатовича.
Все права на вычисления, формулы и тексты принадлежат автору.
Верификация: \texttt{python3 laboratory.py} (пункт <<Стенды>>/<<Сертификаты>>),
Lean: \texttt{verification/lean/HodgeLaboratory.lean}.
}
\end{document}
